% ATTEMPT_STATUS: unresolved
% TOP_PROBLEM: {"problemId": "problem.npt-bound-entanglement-existence", "canonicalTitle": "Existence of NPT Bound Entanglement", "releaseRank": 276, "primaryDomain": "quantum_information_computation", "primaryDomainLabel": "Quantum information & computation", "displayStatus": "open", "releaseStatus": "open"}
% !TeX program = lualatex
\documentclass[11pt]{article}
\usepackage[margin=25mm]{geometry}
\usepackage{fontspec}
\setmainfont{Latin Modern Roman}
\usepackage{amsmath,amssymb,mathtools}
\usepackage{unicode-math}
\setmathfont{Latin Modern Math}
\usepackage{xurl,hyperref}
\hypersetup{colorlinks=true,urlcolor=blue}
\setcounter{secnumdepth}{0}
\setlength{\parindent}{0pt}
\setlength{\parskip}{5pt}
\setlength{\emergencystretch}{3em}
\begin{document}
\section*{\texorpdfstring{Existence of NPT Bound Entanglement}{Existence of NPT Bound Entanglement}}\label{276}
\subsection{Definitions and mathematical statement}
A finite-dimensional bipartite state is a positive semidefinite trace-one operator $\rho$ on $H_A\otimes H_B$. It is entangled if it is not a convex combination of product states. Partial transpose $\rho^{T_B}$ transposes the second subsystem in an orthonormal basis; having a negative eigenvalue is basis independent. LOCC consists of local quantum operations with classical communication. The target asks for $\rho^{T_B}\not\succeq0$ but zero asymptotic rate of distilling maximally entangled pairs from $\rho^{\otimes n}$. Equivalently, in the standard finite-dimensional distillability criterion, seek an NPT state with
\[
 \langle\psi|(\rho^{T_B})^{\otimes n}|\psi\rangle\ge0
\]
for every $n\ge1$ and every vector $\psi$ of Schmidt rank at most two across $A^n:B^n$. This is not merely failure of single-copy distillation.
\par\smallskip\textit{Formulation and concept references:} \href{https://arxiv.org/pdf/2002.03233}{[S1]}.

\subsection{Short English statement}
Can a quantum state have a negative partial transpose yet yield no distillable entanglement, even when arbitrarily many copies and two-way classical communication are available?
\subsection{Research attempt}
\paragraph{Scope and recent-source check.}
GPT-6 Astra Ultra prepared this unresolved attempt on 27 September 2026
using explicit matrix reductions and primary-source inspection. The
target requires non-distillability for every number of copies, across
the grouped bipartition $A^n:B^n$. Source [S1] also discusses a
specific two-copy Werner-state problem. July 2026 manuscripts [E1, E2]
report solutions to that two-copy question in arbitrary dimension.
Their stated scope was inspected; their proofs are not independently
verified or used here. Even accepting a two-copy theorem would not
produce the all-copy state requested by the catalogue. The work below
keeps that quantifier visible and proves elementary restricted cases.

\paragraph{A normalized Werner family and its negative eigenvalue.}
Let $F$ be the swap on $\mathbb C^d\otimes\mathbb C^d$, and for
$0\le\alpha\le1$ set
\[
 \rho_\alpha=\frac{I-\alpha F}{d^2-\alpha d}.
\]
Since $F$ has eigenvalues $+1,-1$, the numerator is positive
semidefinite. Also $\operatorname{Tr}F=d$, so the denominator makes
this a state; for $d\ge2$ it is strictly positive as a scalar.
For $|\Phi_d\rangle=d^{-1/2}\sum_i|i,i\rangle$ and
$P_d=|\Phi_d\rangle\langle\Phi_d|$, direct transposition gives
\[
 F^{T_B}=dP_d,\qquad
 \rho_\alpha^{T_B}=\frac{I-\alpha dP_d}{d^2-\alpha d}.
\]
Its eigenvalue on $\Phi_d$ is negative exactly when $\alpha>1/d$;
all orthogonal directions have positive eigenvalue. This verifies NPT
by an exact spectrum calculation rather than a rounded numerical test.
Because a separable state stays positive under partial transpose, every
one of these NPT states is entangled.

\paragraph{The one-copy Schmidt-rank threshold is exact.}
Represent a normalized bipartite vector by its coefficient matrix $X$,
so $\|X\|_F=1$, its Schmidt rank is $\operatorname{rank}X$, and
\[
 |\langle\Phi_d|\psi\rangle|^2=|\operatorname{Tr}X|^2/d.
\]
For rank at most $r$, the trace is bounded by the sum of the singular
values, which is at most $\sqrt r\|X\|_F$. This proves the sharp
bound $|\langle\Phi_d|\psi\rangle|^2\le r/d$.
Equality is attained by the normalized sum of $r$ matching basis
vectors. For Schmidt rank at most two, therefore,
\[
 \langle\psi|\rho_\alpha^{T_B}|\psi\rangle
 \ge\frac{1-2\alpha}{d^2-\alpha d}.
\]
It follows that the state is one-copy undistillable exactly for
$\alpha\le1/2$, and one-copy distillable for $\alpha>1/2$.
For $d\ge3$, the interval
\[
                         1/d<\alpha\le1/2
\]
contains NPT states which pass the entire one-copy Schmidt-rank-two
test. In dimension two that interval is empty. This is a fully checked
partial result, but passing one copy is not the requested all-copy
conclusion.

\paragraph{The two-copy condition becomes a rank-two matrix inequality.}
For two copies, group the row indices as $A_1A_2$ and the column
indices as $B_1B_2$. Write the coefficient matrix of a vector as
$X\in M_{d^2}(\mathbb C)$, with entries $X_{ia,jb}$.
The relevant restriction is $\operatorname{rank}X\le2$; assume
$\|X\|_F=1$. Define the partial traces algebraically by
\[
 (\operatorname{Tr}_1X)_{ab}=\sum_iX_{ia,ib},\qquad
 (\operatorname{Tr}_2X)_{ij}=\sum_aX_{ia,ja}.
\]
These are partial traces of a coefficient matrix and do not require
$X$ itself to be Hermitian or positive.

Expanding the rank-one maximally entangled projectors gives
\[
 \langle P_{A_1B_1}\otimes I\rangle
                  =\frac1d\|\operatorname{Tr}_1X\|_F^2,
 \qquad
 \langle I\otimes P_{A_2B_2}\rangle
                  =\frac1d\|\operatorname{Tr}_2X\|_F^2,
\]
and the expectation of both projectors is
$|\operatorname{Tr}X|^2/d^2$. Consequently the sign of the two-copy
expectation is the sign of
\begin{equation}\label{a276:twocopy}
 1-\alpha\bigl(\|\operatorname{Tr}_1X\|_F^2+
                    \|\operatorname{Tr}_2X\|_F^2\bigr)
                      +\alpha^2|\operatorname{Tr}X|^2.
\end{equation}
The omitted denominator is $(d^2-\alpha d)^2>0$.
At $\alpha=1/2$, nonnegativity is equivalent to
\[
 \|\operatorname{Tr}_1X\|_F^2+
 \|\operatorname{Tr}_2X\|_F^2
       \le2\|X\|_F^2+\tfrac12|\operatorname{Tr}X|^2
       \quad(\operatorname{rank}X\le2).
\]
This is an exact reformulation, not a proof of the displayed inequality
for every rank-two matrix. It also explains why positivity tests only
on Hermitian or diagonal matrices cannot exhaust the distillability
criterion.

\paragraph{Two restricted tests that can be proved directly.}
First, rank-one vectors across $A^n:B^n$ can never give a negative
expectation for any physical state $\rho$, at any number of copies:
\[
 \langle a\otimes b|(\rho^{T_B})^{\otimes n}|a\otimes b\rangle
 =\langle a\otimes\overline b|\rho^{\otimes n}
                              |a\otimes\overline b\rangle\ge0.
\]
The complex conjugation is in the same basis used for the partial
transpose. This is block positivity; it does not extend from rank one
to rank two by linearity, since the cross terms then matter.

Second, restrict the test vector to factor across copies. Its coefficient
matrix has the form $X=X_1\otimes\cdots\otimes X_n$. The rank is
$\prod_i\operatorname{rank}X_i$. If this is at most two, every factor
has rank at most two and all but possibly one have rank one.
For normalized factors and $\alpha\le1/2$, each single-copy expectation
of $I-\alpha dP_d$ is nonnegative by the preceding threshold proof.
Their product is nonnegative. Thus no copy-factorized eligible vector
can distill a Werner state in the candidate interval, regardless of
how many copies are supplied.

The restriction is strong. A rank-two matrix on the grouped spaces
need not be a tensor product of single-copy matrices. Conversely,
tensoring $n$ arbitrary single-copy rank-two witnesses generally gives
rank $2^n$, which is not an allowed witness in the criterion. Both
mistakes can lead to invalid inductions on copy number.

For a simple genuinely correlated rank-two test, take the diagonal
matrix with entries $1/\sqrt2$ at $(00,00)$ and $(11,11)$ and zeros
elsewhere, for $d\ge2$. Both partial traces have squared norm one,
and $|\operatorname{Tr}X|^2=2$. Formula \eqref{a276:twocopy} becomes
$1-2\alpha+2\alpha^2$, always positive. This verifies one family of
correlated tests without pretending that it exhausts the compact
rank-two variety.

\paragraph{Why tensor positivity is not a valid induction.}
At an NPT parameter the operator $I-\alpha dP_d$ is not positive.
Its negative eigenvector has Schmidt rank $d$, so the one-copy rank-two
restriction hides that eigenvalue when $d>2$. Tensoring operators that
are nonnegative only on a restricted class of vectors does not in
general preserve that restricted positivity: the allowed vectors in
the grouped tensor product have new cross-copy correlations. An
induction would need a decomposition theorem expressing every eligible
global vector through eligible local tests with nonnegative weights.
No such theorem has been established here. Expanding a vector in a
basis gives interference terms, not a convex combination of product
expectations.

Similarly, the presence of a negative tensor-product eigenvalue says
nothing unless a negative vector of Schmidt rank at most two is
exhibited. Searching all eigenvectors of the full operator addresses
ordinary positivity, a different question that NPT has already answered.

\paragraph{Compactness exposes the remaining quantifier issue.}
Let $B_n\subset[1/d,1/2]$ be the set of parameters that pass every
Schmidt-rank-two test at copy number $n$. It is closed: each fixed
vector's expectation is continuous in $\alpha$, and the set is the
intersection of their closed nonnegativity sets. Also
$B_{n+1}\subseteq B_n$. To see this for the Werner family, tensor any
negative $n$-copy witness with a product vector orthogonal to
$\Phi_d$ in the new copy. The extra numerator expectation is one,
so negativity and Schmidt rank at most two are preserved.

If there were a fixed $\delta>0$ such that every
$B_n\cap[1/d+\delta,1/2]$ were nonempty, compactness and nesting would
give a parameter in their intersection. Its state would be NPT and
undistillable at every copy number, answering the problem. The uniform
margin above $1/d$ is essential. Nonempty sets of NPT candidates for
each finite $n$ could approach the PPT endpoint; the abstract nested
intervals $[1/d,1/d+1/n]$ show the logical failure. Their intersection
contains no NPT parameter, despite positive-width NPT portions at
every finite stage.

\paragraph{Verification, gap, and outcome.}
The partial-transpose spectrum, Schmidt-rank bound and coefficient-matrix
contractions were checked independently, including nonsymmetric complex
rank-two examples for the contraction identities. Finite numerical
searches verify only individual vectors and cannot prove universal
positivity on the rank-two variety. Recent two-copy claims do not remove
the all-copy quantifier or establish the uniform-margin condition.
The first remaining gap is an all-copy restricted-positivity mechanism
for one fixed NPT state, or a proof every NPT state eventually fails it.
The one-copy threshold and the stated restricted multi-copy tests are
proved; the existence of NPT bound entanglement remains
\textbf{UNRESOLVED}.

\subsection{Sources}
\begin{itemize}
\item[S1]P. Horodecki, L. Rudnicki, and K. Zyczkowski, Five Open Problems in Quantum Information Theory, PRX Quantum 3 (2022).
\url{https://arxiv.org/pdf/2002.03233}.
\textit{Formulation source.}
\item[E1]Jinshi Fu, Li Gao and Sang-Jun Park, \emph{A solution to 2-copy distillability of Werner states}, arXiv:2607.21367, submitted 23 July 2026.
\url{https://arxiv.org/abs/2607.21367}.
\textit{Abstract and scope inspected 27 September 2026; claimed two-copy result, not used as a proved all-copy assertion.}
\item[E2]Zhiwei Song and Lin Chen, \emph{A partial-trace matrix inequality and Werner-state distillability}, arXiv:2607.23416, submitted 26 July 2026.
\url{https://arxiv.org/abs/2607.23416}.
\textit{Abstract and scope inspected 27 September 2026; proof not independently certified in this attempt.}
\end{itemize}
\end{document}
